% Reader question: If one principal can fund many fake identities, how does
%   the mechanism keep a defaulting identity loss bounded?
% Claim: replicating identities multiplies entry cost, never the loss the
%   mechanism attaches to one default, because slash equals the lesser of the
%   posted deposit and the transaction coverage ceiling.
% Evidence: K interchangeable identities, each depositing B_dep; the closed
%   form s = min(B_dep, B_T) from this chapter section on Sybil deterrence
%   [internal, closed form].
% Must distinguish: entry cost (which rises with K) from attacker loss
%   (which does not rise past B_T no matter how many identities enter).
% Grammar chosen: identity fan-out beside its closed-form response curve.
%   Rejected: the atlas row prescribes an actor sequence with an accounting
%   rail; this page argues a bound on a formula, not an ordered sequence of
%   attacker actions, so the sequence grammar would misstate what this
%   fragment proves. Flagged for the author.
% Five-second test: a reader says slash stops rising once the deposit passes
%   the transaction coverage ceiling, no matter how many identities exist.
\begin{figure}[H]
\centering
\pdfigurehue{pdgold}%
\begin{tikzpicture}[pd figure,x=1cm,y=1cm]
  % ---- one grid -------------------------------------------------------------
  \def\px{0.85}  \def\py{1.35}
  \def\ix{2.75}
  \def\ax{3.55}  \def\bx{4.15}
  \def\ox{4.75}  \def\oy{0.15}
  \def\rx{7.35}  \def\xmax{7.75}  \def\ymax{2.40}
  \def\kx{6.45}  \def\ky{1.85}
  \node[pd title,anchor=north west,text width=6.0cm] at (0,4.45)
    {replicating identities multiplies entry, not one default loss};
  % ---- panel A: one principal fans out into K identities -------------------
  \node[pd identity state,circle,minimum size=8mm,inner sep=1pt] (pp) at (\px,\py) {$P$};
  \foreach \y/\lab in {2.55/1,1.95/2,1.35/3,0.75/{$\ldots$},0.15/{$K$}} {
    \node[pd identity state,circle,minimum size=6.5mm,inner sep=.5pt] at (\ix,\y) {\lab};
    \draw[pd identity rule,dashed] (pp.east) -- ({\ix-.42},\y);}
  \node[pd label,anchor=south] at (\ix,2.90) {$K$ insurer identities};
  \node[pd note,anchor=north,text width=3.6cm,align=center] at (\ix,-.46)
    {$K \times B_{\mathrm{dep}}$ posted; each identity can re-enter};
  \draw[pd focus arrow] (\ax,\py) -- (4.25,\py);
  % ---- panel B: the closed-form response, s = min(B_dep, B_T) --------------
  \draw[pd arrow] (\ox,\oy) -- (\xmax,\oy);
  \draw[pd arrow] (\ox,\oy) -- (\ox,\ymax);
  \node[pd label,anchor=west] at (\xmax,\oy) {deposit};
  \node[pd label,anchor=west] at ({\ox+.08},{\ymax+.16}) {slash $s$};
  \draw[pd guide] (\kx,\oy) -- (\kx,\ky);
  \draw[pd guide] (\ox,\ky) -- (\kx,\ky);
  \node[pd label,anchor=north] at (\kx,{\oy-.10}) {$B_T$};
  \node[pd label,anchor=east] at ({\ox-.10},\ky) {$B_T$};
  \path[pd neutral fill] (\kx,\oy) rectangle (\rx,\ky);
  \draw[pd focus rule] (\ox,\oy) -- (\kx,\ky);
  \draw[pd focus rule] (\kx,\ky) -- (\rx,\ky);
  \node[pd focus datum] at (\kx,\ky) {};
  \node[pd label,anchor=south] at ({(\kx+\rx)/2},{\ky+.14}) {$s=B_T$};
  \node[pd note,anchor=north,text width=3.5cm,align=center] at ({(\ox+\rx)/2},{\oy-.72})
    {extra deposit returns intact past $B_T$};
\end{tikzpicture}
\caption{More identities multiply entry cost, not loss per default.
Slashing stops at coverage $B_T$; deposits above that cap cannot increase
the loss. Figure~\ref{fig:sybil-mc-inline} tests the resulting incentives.
[internal, closed form]}
\label{fig:sybil-inline}
\end{figure}
